%  LaTeX support: latex@mdpi.com 
%  For support, please attach all files needed for compiling as well as the log file, and specify your operating system, LaTeX version, and LaTeX editor.

%=================================================================
\documentclass[geomatics,article,submit,pdftex,moreauthors]{Definitions/mdpi} 

% MDPI internal commands
\firstpage{1} 
\makeatletter 
\setcounter{page}{\@firstpage} 
\makeatother
\pubvolume{1}
\issuenum{1}
\articlenumber{0}
\pubyear{2023}
\copyrightyear{2023}
%\externaleditor{Academic Editor: Firstname Lastname}
\datereceived{11 July 2023} 
\daterevised{8 August 2023} 
\dateaccepted{} 
\datepublished{} 
\hreflink{https://doi.org/} % If needed use \linebreak
%\doinum{}

%=================================================================
% Add packages and commands here. The following packages are loaded in our class file: fontenc, inputenc, calc, indentfirst, fancyhdr, graphicx, epstopdf, lastpage, ifthen, lineno, float, amsmath, setspace, enumitem, mathpazo, booktabs, titlesec, etoolbox, tabto, xcolor, soul, multirow, microtype, tikz, totcount, changepage, attrib, upgreek, cleveref, amsthm, hyphenat, natbib, hyperref, footmisc, url, geometry, newfloat, caption

\graphicspath{{figures/}} % path to folder with figures
\usepackage{subcaption}
\usepackage{listings}
\usepackage{xcolor}
\usepackage{gensymb} % for \textdegree
\usepackage{tabularx}

\usepackage{listings}
\usepackage{xcolor}

\definecolor{deepblue}{rgb}{0,0,0.5}
\definecolor{deepred}{rgb}{0.6,0,0}
\definecolor{deepmagenta}{RGB}{195,12,214}
\definecolor{deepgreen}{RGB}{29,122,30}
\definecolor{mintgreen}{RGB}{192,249,192}
\definecolor{codepurple}{RGB}{204, 0, 153}
\definecolor{LightCyan}{rgb}{0.88,1,1}
\definecolor{GainsboroPurple}{RGB}{247,176,247}
\definecolor{LightSlateGrey}{RGB}{124,118,118}
%    
\lstdefinestyle{mystyle}{
    commentstyle=\color{LightSlateGrey},
    keywordstyle=\color{deepgreen}\bfseries,
    emphstyle=\color{deepred},
    numberstyle=\tiny\color{deepmagenta},
    stringstyle=\color{codepurple},
    basicstyle=\ttfamily\footnotesize,
    breakatwhitespace=false,         
    breaklines=true,                 
    captionpos=b,                    
    keepspaces=true,                 
    numbers=left,                    
    numbersep=5pt,                  
    showspaces=false,                
    showstringspaces=false,
    showtabs=false,                  
    tabsize=2
}
\lstset{style=mystyle}

%=================================================================
% Full title of the paper (Capitalized)
\Title{Seafloor and Ocean Crust Structure of the Kerguelen Plateau from Marine Geophysical and Satellite Altimetry Datasets}
% Interplay Between Seafloor Spreading, Marine Gravity Roughness and Ocean Crust Structure of the Kerguelen Plateau Analysed from Satellite Altimetry

% MDPI internal command: Title for citation in the left column
\TitleCitation{Seafloor and Ocean Crust Structure of the Kerguelen Plateau from Marine Geophysical and Satellite Altimetry Datasets}

% Author Orchid ID: enter ID or remove command
\newcommand{\orcidauthorA}{0000-0002-5759-1089} % Polina Add \orcidA{} behind the author's name

% Authors, for the paper (add full first names)
\Author{Polina Lemenkova $^{1}$*\orcidA{}}

%\longauthorlist{yes}

% MDPI internal command: Authors, for metadata in PDF
\AuthorNames{Polina Lemenkova}

% MDPI internal command: Authors, for citation in the left column
\AuthorCitation{Lemenkova, P.}
% If this is a Chicago style journal: Lastname, Firstname, Firstname Lastname, and Firstname Lastname.

% Affiliations / Addresses (Add [1] after \address if there is only one affiliation.)
\address{%
$^{1}$ \quad Laboratory of Image Synthesis and Analysis (LISA), École polytechnique de Bruxelles (Brussels Faculty of Engineering), Université Libre de Bruxelles (ULB). Building L, Campus du Solbosch, ULB -- LISA CP165/57, Avenue Franklin D. Roosevelt 50, Brussels 1050, Belgium.
}

% Contact information of the corresponding author
\corres{Correspondence: polina.lemenkova@ulb.be; Tel.: +32 471 86 04 59 (P.L.)}

% Current address and/or shared authorship
%\firstnote{Current address: Affiliation 3.} 

% Abstract (Do not insert blank lines, i.e. \\) 
\abstract{The volcanic Kerguelen Islands are formed on one of the world's largest submarine plateaus. Located in the remote segment of the southern Indian Ocean close to Antarctica, the Kerguelen Plateau is notable for complex tectonic origin and geologic formation related to the Cretaceous history of the continents. This is reflected in varying age of the oceanic crust adjacent to the plateau and highly heterogeneous bathymetry of the Kerguelen Plateau with seafloor structure differing for the southern and northern segments. Remote sensing data derived from the marine gravity and satellite radar altimetry surveys serve as an important source of information for mapping complex seafloor features. This study incorporates the geospatial information from NOAA, EMAG2 and WDMAM, ETOPO1 and EGM96 datasets to refine the extent and distribution of the extracted seafloor features. The cartographic joint analysis of topography, magnetic anomalies, tectonic and gravity grids is based on the integrated mapping performed using Generic Mapping Tools (GMT) programming suite. Mapping of the submerged features (Broken Ridge, Crozet Islands, seafloor fabric, orientation and frequency of magnetic anomalies) enabled to analyse the their correspondence with free-air gravity and magnetic anomalies, geodynamic setting and seabed structure in south-west Indian Ocean. The results show that integrating the datasets using advanced cartographic scripting language improves identification and visualization of the seabed objects. The results include 11 new maps on the region covering the Kerguelen Plateau and southwest Indian Ocean. This study contributes to increasing the knowledge of the seafloor structure in the French Southern and Antarctic Lands.
}

\keyword{Antarctic; Southern Ocean; bathymetry; French Southern and Antarctic Lands; cartography; satellite altimetry; marine geophysics; sediments; magnetic anomalies; seafloor} 

\PACS{91.10.Da; 91.10.Op; 93.85.Pq; 91.10.Fc; 91.10.Jf; 91.10.Vr}
\MSC{86A30; 86-05; 86A04; 86A70}
\JEL{Y91; C83; C88; C90; C00; C60; C61}

\begin{document}

%----------- section ----------->
\section{Introduction}

%----------- subsection ----------->
\subsection{Background}

\hl{In cartography, spatial data processing, plotting maps (also mapping layouts) are widely used as a common task, i.e. visualization and representation of spatially defined objects using cartographic techniques} \cite{PAPPVARY1989103,ROBINSON198991,GOODCHILD20091}. \hl{One of the generally accepted methods of mapping and quantitative and qualitative cartographic visualization is implemented via the cartographic software named Geographic Information Systems (GIS) using algorithms of raster and vector data processing embedded in these programs} \cite{GOODCHILD2009500,SLOCUM1991167,TAYLOR1994333}. \hl{Obviously, the high complexity and time-consuming performance of GIS are not conductive for their large-scale applications of mapping in various tasks of Earth sciences. In order to reduce the data processing complexity in computer graphics and time cost of mapping through automation, programming and scripting algorithms are proposed as an alternative to the GIS-based approach to mapping} \cite{10162166}.

\hl{The application of programming methods in cartography has been extensively developed and resulted in released several programs utilizing scripts as a tool of mapping. Programming-based cartographic approaches can be divided into three general types. First, partially using scripts (e.g., as plugins or alternative tools additional to the existing Graphical User's Interface (GUI) in classical mode). These include, for instance, ModelBuilder by ESRI} \cite{6344536}, \hl{modules of Geographic Resources Analysis Support System GIS (GRASS GIS), processing script editor in QGIS or ArcGIS} \cite{OMRAN2016140,NAGHIBI2021101232}\hl{, or the Python-based ERDAS Macro Language (EML) for scripting used in Erdas Imagine which enables creating spatial models. Second, spatial libraries of programming languages. For instance, these include selected packages of R and Python specially designed for satellite image processing and geospatial data analysis} \cite{ROBERTS2022105192,jimaging8120317,VOS2019104528,app122412554}. \hl{Third, the programs that are completely based on using scripting language without any graphical user's interface (GUI) such as Generic Mapping Tools (GMT). All these examples of using scripts aim at automating the cartographic data processing by advanced scripting tools.}

\hl{The script-based cartographic programs are based on the principle of using syntax of programming language which includes a number of key commands and expressions recognisable by the system} \cite{6057428}. \hl{Using the rules of the embedded language, it becomes possible to write the script for data modelling and cartographic visualization. In contrast to the traditional GUI-based GIS methods, cartographic scripts do not plot the maps directly, however, their provides a series of commands which include essential information on appearance of maps which controls the specific features on a plotted cartographic layout} \cite{Lemenkova202215}. \hl{In fact, scripts and programming commands used in cartography define the elements on map produced by script during execution. Specifically, it is possible to define the key map concepts such as mathematical definitions of the map (projections, resolution, grid, coordinate systems extent and scale), design of symbols and legends (colours and size of the objects, palettes for continuous fields and transparency), and exposition of the elements (overlay, topology, generalisation)} \cite{ABER201753}.

\subsection{\textcolor{blue}{Problem formulation}}
The analysis of the seafloor structure and ocean crust concerns mapping, detecting, and recognising the bathymetric structures and variations in geophysical fields\hl{. }Among these problems, the integrated geophysical interpretation of the satellite altimetry, marine gravimetry and magnetic anomalies as well as acoustic and seismic surveys present one of the most active research areas which have attracted\hl{ }research interest over the past decades \cite{Bascou}. Previous studies \cite{ESHAGH2021313} provided a broad overview of\hl{ }approaches for analysis of the satellite gravimetry data to model the Earth beneath the sea. Marine geophysical and satellite altimetry data provide information on crustal density and processes in the upper mantle \cite{KAPAWAR2021106692,PENG2022110823,ZHANG202368}, sediment thickness and basement depth \cite{Lemenkova202005}. Moreover, the analysis of the geophysical data enables to reveal key characteristics regarding the lithosphere thickness such as Moho discontinuity and elastic features of the lithosphere \cite{TESAURO201378}, viscosity and flexural rigidity. Such data are essential for\hl{ }investigations \hl{on} the Earth's structure.

In terms of the data-driven analyses, previous works investigating the oceanic seabed can be roughly classified into studies focused on the bathymetric mapping and geophysical analysis. Mapping seafloor\hl{ }characterises the bathymetric patterns \hl{using }data visualization and\hl{ }cartographic methods \cite{Battista,Lemenkova202105,Lemenkova202012}. Marine geophysical methods\hl{ }investigate the geology of the continental margins using methods of the deep-sea drilling \cite{Schlich1992TheGA,Veevers1991a}, the analysis of features extracted from seismic survey and remote sensing data \cite{NEGI2022229330}, investigating the structure of the deep ocean basins \cite{BENARD2010633,RAMANA2001241,DAVIS201616} and modelling the mid-ocean ridge system\hl{ } \cite{FREY200073,Sreejith,Lemenkova202112}. Although \hl{traditional methods of bathymetric survey } supply novel information on origin, evolution, structure, stratigraphy, and tectonic features of the oceanic crust, they\hl{ }require \hl{costly} hydrographic equipment such as mutlibeam echo sounder systems\hl{ }\cite{Falvey}.

%----------- subsection ----------->
\subsection{Related work}

\hl{Seafloor} bathymetry and geophysical setting of the oceanic crust are two important topics for modelling lithosphere \cite{WATSON201697}. However, the integrated capture and processing of the multi-source data, -- such as bathymetry, marine geologic data (development on the Earth's crust including age, spreading rates and symmetry) and marine geophysical data (sediment thickness, gravity and magnetic anomalies), -- has always been a challenging task, and few attempts have been made to explore it. \hl{Since the variety of data sources is increasing in modern cartography, it results in an exponential growth in the volume of data serving cartographic applications} \cite{LEE199121}\hl{. Therefore, }the use of the remote sensing data, satellite altimetry and gravimetry has proven to be effective in geophysical\hl{ }studies. Numerous examples exist in\hl{ }literature \hl{that focus on} estimating the geoid values, evaluating gravity anomalies, analysis of the tectonic and crustal structures, glacial and hydrological modelling and habitat mapping. 

Satellite altimetry data can be used to jointly represent the patterns of the oceanic currents and to draw the conclusions on\hl{ }connectivity between the habitats \cite{MORI201645}. For instance, a study by \cite{COTTE2022103650} \hl{ uses} the hydrographic and acoustic data\hl{ }for the analysis of the vertical layers of the ocean to identify the community distribution\hl{. }An example of the tide modelling \hl{of} bathymetric gradients was presented using the satellite altimetry data \cite{Maraldi}. The retreat of glaciers using mass balance measurements was estimated using remote sensing data to analyse spatial distribution of the surface mass balance \cite{Verfaillie}. Other studies used the \hl{Gravity Recovery and Climate Experiment} (GRACE) mission for detecting trends and variations in the ice-sheet mass balance by evaluating gravity signals \cite{YAMAMOTO2008267}. Furthermore, Mathieu et al. \cite{MATHIEU2011206} combine the SRTM DEM, remote sensing data and\hl{ }geological \hl{sampling} to clarify \hl{seafloor} structures\hl{.} 

These and similar examples illustrate that a combination of the bathymetric and geophysical data for detailed analysis of the Earth's crust structure and topography outperforms their use separately, and suggests that seafloor mapping should be based on the complex analysis of the heterogeneous features of the seabed. Such information can be derived from various datasets on Earth \hl{observation}, as pointed earlier \cite{COWLEY2004267,Lemenkova202125,1993A234,ABUAMARAH2019868}. Since such methods can be regarded as feature-level combinations of the seafloor fabric showing local appearance of the prominent features of the seabed such as fracture zones, ridges, seamounts, deep-sea trenches, canyons \hl{or} rifts,\hl{ }the use of\hl{ }programming approach to merge different integrated geophysical datasets enables to match \hl{multi-source} data\hl{ with }different resolution and origin. Such data can be used for a detailed analysis of the seafloor structure and geodynamics \cite{Veevers1991b,Lemenkova202104,Hill}. 

%----------- subsection ----------->
\subsection{Objection and motivation}

In this study, several bathymetric and marine geophysical datasets were \hl{used} for the analysis of the Kerguelen Plateau (Figure \ref{fig01})\hl{ }within the south-western segment of the Indian Ocean to visualise, describe and analyse the structures of the seafloor\hl{ }in the context of the geophysical settings.\hl{ }Ten new maps are presented and described to visualise spatial variations, reveal correlation and\hl{ }matches between the geophysical and topographic setting of the Kerguelen region to\hl{ }continue the existing \hl{similar studies} \cite{Nougier,Petit,RECQ1989133}. Given the remote location of the Kerguelen Archipelago, one of the most isolated places on Earth, and associated difficulties and high cost of geological and geophysical sampling, the integrated use of the satellite derived data enables to get better insights into the structure of the Kerguelen Plateau\hl{.}

%----------- section ----------->
\section{Study area}

The Kerguelen Plateau presents a large topographic elevation \hl{in south-west Indian Ocean }which extends on $6500 km^2$ \cite{Operto}\hl{, Figure }\ref{fig01}. The Kerguelen Archipelago formed on the plateau forms a shallow submarine plateau. The structure of the Kerguelen Archipelago is asymmetric with notable variations in the two distinct parts: northern part (Heard, McDonald and Kerguelen Islands) is more shallow (<1000 m) compared to the southern segment (depths of 1500-2000 m) \cite{Ramsay}.\hl{ }Morphologically, the \hl{plateau} is oriented in \hl{a} NNW direction toward the Antarctic continental margin where it is constrained by the Princess Elizabeth Land. On the northwest, the archipelago is limited by the Crozet Archipelago which is located 1250 km \cite{BRETON2013126}, to the north -- by the African-Antarctic Basins, to the northeast -- by the Australian-Antarctic Basin, and to the south -- by the Antarctic. \hl{Adjacent} lands include the Enderby Land, Kemp Land, Mac Robertson Land, Princess Elizabeth Land, Wilhelm II Land and Queen Mary Land \hl{located in }the Antarctic, Figure \ref{fig02}. 

\begin{figure}[H]
\centering
	\hspace{100pt}\includegraphics[width=14.0 cm]{Fig_01.jpg}
\caption{Map of the Kerguelen Plateau region. Mapping: GMT. Map source: author. \label{fig01}}
\end{figure}

The origin of the Kerguelen Plateau is related to the hotspot associated with the\hl{ }Gondwana breakup in Cretaceous \hl{(}120--110 Ma)\hl{, }and \hl{seafloor }spreading between India and Antarctica in\hl{ }south-western\hl{ }Indian Ocean \cite{Guglielmi}. Early surface expression of the Kerguelen Plume \hl{in} the south\hl{ }initiated in Mesozoic \hl{in} 120 Ma \cite{Gaina}. Active submarine volcanism on the aseismic ridge and associated sedimentation processes since Cretaceous resulted in the formation of the basaltic basement of Kerguelen \cite{Henry}. The volcanism continued in the Kerguelen hotspot up to Tertiary (<40 Ma) and is\hl{ }reflected in\hl{ }geologic and mineralogical structure of the archipelago\hl{: }the archipelago is covered up to 85\% by basalts originated from the mantle plume of the hotspot\hl{, }resulting from the crystallisation of the raised and cooled basaltic magma \cite{NICOLAYSEN2000313,Lacroix,Doucet}. Currently, \hl{the Kerguelen Plateau} presents one of the World's largest volcanic plateaus with a total length of 2300 km. Occasional evidence of volcanism is still remaining on the Heard and McDonald Islands \cite{Li}.  

The Kerguelen Plateau belongs to the French Southern and Antarctic Lands with protected environment \cite{Rue1931,Faivre,Dupouy}. \hl{Its }high environmental value\hl{ }is explained by the unique wildlife structure\hl{ }which includes vulnerable species. The remote location of the archipelago -- 4.000 km from South Africa and Australia  \cite{Rue1932b,fishes8040174} created \hl{perfect }conditions to preserve\hl{ }unique flora and fauna of \hl{the Kerguelen Archipelago}. Rare \hl{Antarctic }plant species and high biodiversity\hl{ are }associated \hl{here }with soil developed on the basaltic basement and specific geologic substrate of volcanic origin\hl{. Other reasons include} the effects from \hl{the }Polar climate\hl{ }\cite{Rue1932a,Paulian,Aubert3723}. The distribution of fish communities and phytoplankton is affected by seasonal changes of the Antarctic circumpolar currents \cite{Park2009,Damerell,MONGIN2008880,Park,Park2014,POLLARD20071915,Pauthenet,Penna}. \hl{Moreover, }steep bathymetric slopes and exposition of the Kerguelen plateau \hl{strengthen the speed and intensify} circulation\hl{ }of the ocean currents \cite{Rietbroek}. As a result, \hl{a} combination of all these factors makes the ecosystems of the Kerguelen Archipelago unique and protected as an environmental heritage of the French Southern and Antarctic Lands \cite{Sombetzkib,Sombetzkia,Naim}.  

\begin{figure}[H]
\centering
	\hspace{100pt}\includegraphics[width=14.0 cm]{Fig_02.jpg}
\caption{Age of the ocean crust in SW Indian Ocean and the Kerguelen Plateau region and east Antarctic. Mapping software: GMT v. 6-1-1. Map source: author. \label{fig02}}
\end{figure} 

%----------- section ----------->
\section{Materials and Methods}

All the maps have been plotted using scripting toolset Generic Mapping Tools version \hl{6.4.0}, \href{https://www.generic-mapping-tools.org/}{https://www.generic-mapping-tools.org/}. A key aspect of the GMT algorithms is that is considers each cartographic element by its parameters and adds new layers on the map regarding the target location which can be adjusted using refined flag options in special modules. Thus, a modular scripting approach of GMT differs it from GIS. 

%----------- subsection ----------->
\subsection{Data}

The materials used\hl{ }as input cartographic grids include the following datasets. The marine geological data on age, spreading rates, and asymmetry of the ocean crust were derived from the NOAA high-resolution dataset \cite{Mueller}, Figure \ref{fig03}. The units of the dataset on spreading asymmetry are expressed as \% of crustal accretion of the seafloor, i.e. symmetric spreading results in values of 50\% on conjugate ridge flanks. The bathymetric mapping was based on the ETOPO Global Relief Model \cite{NOAA}. The gravity grids were obtained from the EGM-2008 and EGM96 \cite{Pavlis}. The gravity data are based on the principle of gravity anomaly over various surfaces of the Earth which is reflected in the relationship between the topography, distribution of masses within the local \hl{and regional }relief \hl{structures }and \hl{marine }gravity values, and the development of gravitational potential in various locations of the Earth. 

\begin{figure}[H]
\centering
	\hspace{100pt}\includegraphics[width=13.5 cm]{Fig_03.jpg}
\caption{Asymmetry in crustal accretion on conjugate ridge flanks of the ocean crust over the Kerguelen Plateau region, east Antarctic and south-west Indian Ocean. Map source: author. \label{fig03}}
\end{figure}

The free-air and Bouguer gravity anomaly grids are derived from \hl{high-resolution grids} \cite{EGU,BRAITENBERG2021706}. Specifically,\hl{ }gravity field differs over the oceans due to variations in density. The highest anomalies arise from undulations on density,\hl{e.g.,} at the\hl{ }heterogeneous seafloor or at the crust--mantle interface (Moho discontinuity) \cite{MCNUTT201920}. Such data enable to model\hl{ }the gravity fields for analysis of the Earth structure.\hl{ }Current updates of gravity grid improve the altimetry-derived gravity anomalies. Although gravity grids do not indicate the topography directly, \hl{they provide} essential information on \hl{the }relief over the Earth which gives an approximate in relationship between \hl{the }topographic and geophysical data and the agreement of surface gravity values with local geoid undulations \cite{CONDIE202281}. 

The data on sediment thickness are collected from the GlobSed dataset which \hl{models }distribution of the sediment thickness in the World's Oceans \cite{Straume}. The magnetic data are derived from the two available information sources compiled originally from \hl{the }satellite and marine magnetic measurements: the Earth Magnetic Anomaly Grid (EMAG2) \cite{Maus} which has a resolution of 2 arc minutes, and the World Digital Magnetic Anomaly Map (WDMAM) \hl{with} 3-minute resolution \cite{Lesur}. The gravity data were derived from the free-air global gravity grids \cite{Sandwell}.

The analysis of the geodetic data is useful in such complex terrains since it enables to identify major tectonic structures and topographic patterns. Besides the EGM grids, the identification of the lithospheric structure\hl{ }and tectonic blocks can also be performed using\hl{ }satellite gravity data such as World Gravity Model (WGM) and\hl{ }GRACE\hl{ }\cite{EMISHAW2023104775}. Furthermore, the EGM and astrogeodetic vertical deflections are useful for modelling gravity and geopotential difference \cite{WANG2023}. Finally, processing the terrestrial and satellite geodetic data can be employed for solving theoretical issues of geodesy such as altimetry–gravimetry boundary-value problem \cite{Lehmann}. Other examples of implications of the geodetic datasets include the use of the satellite-derived data from Cryosat-2 and altimetry with ship-measured gravity\hl{ }for estimating \hl{the }marine gravity anomalies \cite{NGUYEN2020505}.

%----------- subsection ----------->
\subsection{Methods}

In this study, the cartographic methodology is based on using the Generic Mapping Tools (GMT) programming suite version 6.4.0 \cite{Wessel} by the developed scripting workflow explained in detailed in earlier works \cite{Lemenkova202212,Lemenkova202203}. The most prominent feature of the GMT is\hl{ }a scripting approach that principally differs it from the conventional software due to the embedded programming language  \cite{WesselLuis}. The traditional GIS-based methods either employ the\hl{ }GUI\hl{ }for mapping\hl{ }with existing standard\hl{ }menu\hl{ }or allow scripting as a complimentary workflow. 

In contrast, the GMT is a completely console-based software which operates entirely using \hl{scripts}. In this way, \hl{it} captures the \hl{data} by running a script written using \hl{the }embedded \hl{syntax} which operates similar to the programming\hl{. }Script consists \hl{in} a consequence of several commands with pre-defined parameters that control the appearance of the cartographic elements \hl{and features}. The most well-known cartographic tools that use scripts for data processing are GMT \cite{Lemenkova202217} and GRASS GIS \cite{NETELER2012124}\hl{. Other} software\hl{ }allow scripts as an additional functionality, besides the standard \hl{GUI }(e.g., ArcGIS or QGIS). While the latter is mostly used for image processing,\hl{ cartographic} and ecological studies\hl{ }\cite{ROCCHINI2017166,FRIGERI20111265,Lemenkova202313,SOROKINE2007685}, \hl{the }GMT is applied in geophysical and geological studies \cite{Lee,Lemenkova202206,DeSanto}.

\subsubsection{Topographic/bathymetric mapping}

The script for plotting the topographic/bathymetric map is presented in Listing \ref{lst01}. The most essential element of the code are as follows. Mapping starts with selecting the\hl{ study area} from the global ETOPO1 grid by 'grdcut'. The selection of the isolines is performed through a 'grdcontour' with 1.000 m gap. In our case, these lines are meaningful locations that yield a high representation estimate for seafloor bathymetry. The use of 'grdimage' results in plotted image of the raster grid in colours defined by 'makecpt' module\hl{ with the }values -T\hl{ indicating} extreme (min/max) data for topographic grid inspected by 'gdalinfo' module\hl{. }The rest of the code is presented in Listing \ref{lst01} with key comments.

\subsubsection{Mapping seafloor age, spreading rates, spreading asymmetry and age uncertainty of the ocean crust}

Mapping age, age uncertainty, spreading rates and spreading asymmetries of the Kerguelen Plateau\hl{ }is performed based on the raster grids with 2 minute resolution, Figure \ref{fig04}.  The grids are reprojected to Lambert Azimuthal Equal-Area projection using data from the major features in the basin of Indian Ocean. The vector lines of the mid-ocean ridges, tectonic plates and lithospheric plates are added using the 'psxy' module to show the borders of the Indian, African and Australian plates\hl{, }Figure \ref{fig04}. Spreading half rates of the ocean crust are visualised based on the 2 arc min netCDF grid v.3.6 (NOAA). Likewise, the tectonic slab contours are added using the 'psxy' module of the GMT\hl{. }Several grids are combined and represented through plotted raster grids to show the age-depth relationships in the seafloor around the Kerguelen and southwest Indian Ocean. The full scripts are presented in the Appendix \ref{App} of this study with Listing \ref{lst02} used for mapping age of the oceanic lithosphere over Kerguelen Plateau, Listing \ref{lst03} -- for mapping asymmetries in crustal accretion on conjugate ridge flanks over Kerguelen Plateau, Listing \ref{lst04} -- for visualisation of spreading half rates of the oceanic lithosphere in Kerguelen Plateau and SW Indian Ocean\hl{, }and Listing \ref{lst05} -- for mapping crustal age uncertainty of the oceanic lithosphere over the Kerguelen Plateau, Figure \ref{fig05}.

\begin{figure}[H]
\centering
	\hspace{100pt}\includegraphics[width=14.0 cm]{Fig_04.jpg}
\caption{Spreading half rates of the ocean crust over the Kerguelen Plateau region, east Antarctic and south-west Indian Ocean with added GSFML data. Mapping tool: GMT. Map source: author. \label{fig04}}
\end{figure}

\subsubsection{Mapping sediment thickness}

A representation of the GlobSed 5-arc-minute total sediment thickness grid is extracted from the global raster grid on a target location of the Kerguelen Archipelago area and surrounding regions using \hl{'grdcut'} command which selects the file. The amplitude of values is checked using the Geospatial Data Abstraction Library (GDAL) library embedded in GMT \hl{using the} 'gdalinfo'\hl{. }Given this extent of the values, the colour palette table is defined accordingly and presented using the 'psscale' module with defined the parameters of the geographic location for each element on the map\hl{ where} the '-Dg' command defines the position of the colour scale on the map by coordinates. The grid annotations and graticules are added using the 'psbasemap' GMT module. The rest of the script with added comments for brief explanations is presented in Listing \ref{lst06} in the Appendix of this study. This listing was used for plotting the map in Figure \ref{fig06}. 

\subsubsection{Geophysical mapping}

To highlight the features and subsurface characteristics not visible from the geological data, the maps of gravity anomalies were prepared using gravity datasets in the Lambert Equal-Area Azimuthal Projection for conformity with other maps. The full scripts for plotting geophysical maps are presented in Appendix in Listing \ref{lst07} for free-air gravity anomalies and Listing \ref{lst08} for vertical gravity gradient. The most essential lines of the code are as follows. For free-air gravity, the module 'img2grd' was first applied to convert the image from the .IMG into GRD format. The projection was defined using '-JU43/6.0i' flag in 'grdimage' module and the region extent is clipped by \hl{'-R' command} which passed to the next modules accordingly using '-R -J' flags. The vertical gravity gradient was plotted using the satellite-derived grid\hl{ }obtained from the altimeters CryoSat-2, Jason-1/2 \cite{Harper} which enables to identify ridge propagation on the seafloor.

\begin{figure}[H]
\centering
	\hspace{100pt}\includegraphics[width=14.0 cm]{Fig_05.jpg}
\caption{Age uncertainty in lithosphere crust (m.y) over the Kerguelen Plateau region, east Antarctic and south-west Indian Ocean. Cartography: GMT, version 6-1-1. Map source: author. \label{fig05}}
\end{figure}

Since GMT only considers each feature once and processes it using relevant module \hl{(}'psbasemap', 'grdcontour', 'psscale'), its running speed is high. Such technical advantage enables \hl{to process} the grids that cover large geographic areas in a few seconds. With prior knowledge of the geological and geophysical structure of the seafloor, GMT is fast and achieves good mapping performance. \hl{Since} the algorithm adjusts the projection parameters\hl{, it} sets focus on the main seafloor features\hl{ }in the required scale and projection. Furthermore, the GMT-based framework performs well for \hl{defining} the location, extent and borders of the mid-ocean ridges and major basins, \hl{as} the redundant bathymetric features can be removed by cartographic generalisation \hl{using} the gap in the isolines. This enables to find strongly relevant and nonredundant features, highlight\hl{ }correlation between the geophysical setting, gravity and magnetic anomalies, and compare them with geodetic anomalies\hl{. }Hence, \hl{the }seafloor structures can be mapped and visualised\hl{ }using effective GMT approach. 

\subsubsection{Mapping magnetic anomalies}

\hl{To map} the magnetic anomalies, the two different grids were used -- the WDMAM and the EMAG2. The first was subset from the global grid embedded in the GMT using the 'grdcut' \hl{module}, and the second was obtained from the existing raster file of EMAG 2-minute resolution grid\hl{. Hence, } data processing in GMT is flexible and possible using both ways: by handling existing files as stored data, and using embedded grids which are available in the\hl{ }system using remote access to these files. \hl{The} images were \hl{then }visualised using the 'grdimage' module with color palettes adjusted according to the range of data distribution over the Kerguelen Plateau and adjacent south-west region of the Indian Ocean using the \hl{'makecpt' module}. The remaining parts of the GMT scripts are presented for both maps visualised using the WDMAM and the EMAG2\hl{ }in Listing \ref{lst09} and \ref{lst10} in the Appendix.

\subsubsection{Mapping geoid anomalies}

The map of the geoid inundations is completed using the EGM96 dataset using the\hl{ }GMT script where each module defines certain features on a map. To make mapping\hl{ }effective, the scripts are run from the console for \hl{plotting} the geoid and adjusted accordingly to the target \hl{concepts}: projections, grids, colour palettes, raster extent and annotations\hl{. }To this end, these cartographic elements were also defined stepwise on the presented maps of geoid, as \hl{showed} in a script\hl{, see} Listing \ref{lst11}. Such strategy of data visualization\hl{ }ensures high automation of data processing \hl{by GMT }which outperforms the conventional\hl{ }tools.

Strongly relevant features, such as geoid undulation, seafloor isochrons, basement depth isolines, and asymmetries in crustal accretion are associated with the asthenospheric flow from mantle plumes of Kerguelen to spreading Southwest and Southeast Ridges. Mapping such data enables to reveal information about the seafloor attributes and compare it with the geologic structure and geodynamics. Therefore, mapping multi-source topographic and geophysical datasets enables to compare such features using comparable projection. In reality, however, seafloor features are rarely identical due to many factors affecting the bathymetry but may be strongly correlated. To this end, the GMT scripts can be adjusted to other scale and map the features in the enlarged view for detailed order of visualisation. To evaluate the impact of geophysical and geological setting on seafloor structure using the performance of GMT scripts, a number of cartographic trials and re-projections was generated in which each map represented the selected seafloor features with various level of details, colour table for tuned visualisation. The optimal projection was selected as Lambert azimuthal equal-area, which was applied for all the maps for comparability and conformity.

%----------- section ----------->
\section{Results and Discussion}

Automatic matching of topographic and geophysical grids with high accuracy is essential for complex geologic-tectonic investigations. This paper demonstrated the use of scripting algorithms of GMT which were applied for plotting 10 thematic maps covering the south-west segment of Indian Ocean and East Antarctic using bathymetric, geodynamic and high-resolution geophysical datasets on gravity (Figure \ref{fig07} and \ref{fig08}). 

These maps reveal the structure of the seafloor in the Kerguelen Plateau and the surroundings region. The results confirm that the dynamics of the seafloor development reflected in the maps of oceanic crust age, asymmetries, spreading rates of the south-west and southeast Indian ridges are related to major geophysical setting as reflected on the topographic, magnetic and gravity grids. Furthermore, the volcanic activity of the Kerguelen hotspot makes the impact on the distribution of the magnetic anomalies, which supports relevant previous studies \cite{Fanjat,Weissel,Jokat}. Overall, the results demonstrate that the GMT scripting is a powerful and stable cartographic method that performs efficiently for geophysical and bathymetric seafloor mapping. In the sections below, the obtained results are discussed on relevant maps with provided comments on the essential features and traits of the seafloor around the Kerguelen Plateau, south-west Indian Ocean and the East Antarctic. 

%----------- subsection ----------->
\subsection{Ocean floor formation}

The general physical-geographic structure of the Kerguelen Plateau is visible in the map of Figure \ref{fig01}. It shows a morphological orientation of the archipelago in the NW-SE direction and a total extent surpassing 2000 km in length. Northern and southern parts of the plateau are asymmetric where less expressed southern part is older and lie in deeper water in the topographically downlifted areas. Age, spreading rates and spreading symmetry of the ocean crust indicate the evolution steps of the ocean floor formation in southwest Indian Ocean and around the Kerguelen Plateau. The geodynamics setting of the oceanic crust (Figures \ref{fig02}, \ref{fig03}, \ref{fig04} and \ref{fig05}) shows a strong relationship between the Kerguelen Plateau and the two mid-ocean ridges. 

The age of the ocean crust was determined by the interpolation of the adjacent seafloor isochrons oriented towards the direction of seafloor spreading, Figure \ref{fig02}. Such correspondence highlights that unique geophysical setting of the Kerguelen Plateau underlain by the oceanic crust is strongly associated with its tectonic origin associated with volcanic hotspot and geologic history. The formation of the Southwest and the Southeast Indian Ridges is related to the uplift of the Kerguelen Plateau as a remnant of the Mesozoic oceanic basin existing after the separation of Gondwana. The analysis of the bathymetric map with geodynamic maps shows that seafloor heterogeneity around the Kerguelen Archipelago correlates with the seafloor spreading rates where rougher basement is formed in the areas with the low half-spreading rate threshold (30-35 mm/year and lower) which can be revealed by comparison of maps in Figure \ref{fig01} and Figure \ref{fig04}. Thus, the heterogeneity in seafloor patterns varies significantly in various basins of the Indian Ocean depending on the geodynamic setting and the geologic development of the ocean crust.

Gridded map of age of the ocean crust around the Kerguelen Plateau (Figure \ref{fig02}) shows a correlation with the observed spreading half rates of the lithosphere (Figure \ref{fig04}). Originally formed as a single structure, Kerguelen was then split by the seafloor spreading in the south-west sector of the Indian Ocean which resulted in the formation of the two segments of the archipelago \cite{BRADLEY1988243}. Moreover, the difference in the volcanic activity on the northern and central parts of the Kerguelen Plateau underlying the Heard Island proves that it is located on a hotspot with various parts of the islands experiencing the effects from the mantle plume in different scale \cite{Duncan2016,Plenier,HOUTZ197795}. Such unstable position impacts the geodynamic patterns over the Kerguelen Islands, with higher uncertainty of age of the oceanic crust, Figure \ref{fig05} compared to the adjacent areas. Further, the structure of the oceanic crust beneath the Kerguelen Plateau is similar to the crust beneath the aseismic ridges such as the Crozet Rise and the Madagascar Ridge, which proves that it originates by active volcanism from a hot spot, as noted earlier \cite{Recq}.

%----------- subsection ----------->
\subsection{Sediment thickness}

Mapping sediment thickness using the 5-arc-minute GlobSed grid relies on the approximated modelling which highlights the sedimentation trends around the Kerguelen Archipelago, Figure \ref{fig06}. The distribution of the sediment thickness correlates with the age of the underlying oceanic lithosphere and its latitude, which can be noted by the comparison of maps in Figure \ref{fig02} and Figure \ref{fig06}. Such correspondence is especially visible for higher values of the sediment thickness near the shorelines of Antarctic, Amery Ice Shelf (3000-4000 m), Enderby Land (over 4000 m) and the Kerguelen Plateau (2000-3000 m). Higher level of sediment thickness in these areas may also indicate earlier processes of the subaerial erosion which took place before the subsidence and associated sedimentation. 

Rifting during the Late Paleogene resulted in changes in the Tertiary sediment facies of Kerguelen which are affected by the evolution of the Antarctic environment \cite{Froehlich}. Sediments covering the Kerguelen Plateau include pillow lavas, tuffaceous sediments and marine siltstones which were deposited since late Miocene \cite{Duncan}. They are continued as a thick sequence of the Cenozoic sediments (over 5000 m) within the Enderby Basin to the south-west off the Kerguelen Plateau, Figure \ref{fig06}. The correlations reported between the sediment thickness (Figure \ref{fig06}) and the age of the oceanic lithosphere (Figure \ref{fig02}) demonstrate the role of the ocean floor formation on the accumulation of the sediments. 

\begin{figure}[H]
\centering
	\hspace{100pt}\includegraphics[width=14.0 cm]{Fig_06.jpg}
\caption{Sediment thickness (m) in East Antarctica, South-West Indian Ocean and the Kerguelen Plateau region. Cartography: GMT. Lambert azimuthal equal-area projection. Map source: author. \label{fig06}}
\end{figure} 

Other important factors increasing the accumulation of the sediments around the Kerguelen and the adjacent area include glacial processes and the turbidity of the ocean currents. Hence, intense circulation results in the accumulation of the large sediment fields with values over 3000 m. This results from the developed Antarctic circumpolar currents which started initially around the Eocene--Oligocene and continued until present time as the strongest current system in the oceans directed clockwise around the South Pole affecting the adjacent sub-Antarctic regions such as Kerguelen \cite{Huang,ROBERT200237}. A general orientation of maximal sedimentation areas is coherent to the sediment-filled troughs stretching in NW-SE-direction which are associated with general NW-SE orientation of Kerguelen and axes of the tectonic faults. High values in sediment thickness around the Kerguelen Plateau controured by the ridge isolines its eastern margin is associated with the depositions from the bottom currents directed westwards.

%----------- subsection ----------->
\subsection{Free-air gravity anomaly and vertical gravity gradients}

The visualised marine gravity field over the Kerguelen Plateau region and the adjacent areas of the south-west Indian Ocean are shown in Figure \ref{fig07}. The comparison of the gravity roughness with the map of the half-spreading rates (Figure \ref{fig04}) and sediment thickness (Figure \ref{fig06}) shows the relationship between the speed of the spreading of mid-ocean ridges and roughness of the seafloor basement which is explained by the process of mid-ocean ridge formation. Other factors include the associated magma flows, spreading directions in Mesozoic and isochron orientations of the age of the oceanic crust which affect current bathymetric and gravity patterns \cite{Whittaker}. Furthermore, the gravity highs around the Kerguelen Plateau and the Heard Island correspond to the maximal bathymetric elevations and indicate on the seamounts formed of Miocene basalts erupted during the volcanic activity in the southern Indian Ocean and resulted in a large igneous province \cite{Weis}.

\begin{figure}[H]
\centering
	\hspace{100pt}\includegraphics[width=14.0 cm]{Fig_07.jpg}
\caption{Marine gravity field over the Kerguelen Plateau region, south-west Indian Ocean. Mapping: GMT. Map source: author. \label{fig07}}
\end{figure}  

The analysis of maps shows that the values of the free-air gravity over the Kerguelen are higher that in the surrounding areas and reach up to 80 mGal (dark orange colour, Figure \ref{fig08}). In contrast, lower values are associated with the bathymetric depressions and have values of -40 to -60 mGal. Furthermore, the abyssal plain is characterised by the medium values of 0-20 mGal, Figure \ref{fig08}. This well illustrates the correlation of the free-air gravity anomalies with the distribution of topographic highs and depressions on the seafloor since they are strongly influenced by a gravitational effect of the distributed topographic masses since they are caused by differences in elevation. 

Figure \ref{fig08} shows the outputs of the mapping for vertical gradient over the Kerguelen and the southwest Indian Ocean where the effects of different locations are presented. The visualised map demonstrate that the crucial property of gravitational systems like free-air gravity is not only the inclusion of geographic location and the latitude of the selected measured regions, but also the altitude of the Earth's surface since greater altitude implies greater distance from the Earth's centre with affects gravity values. Vertical gravity gradient identifies the variations in gravity with the change in topographic elevations are shown in Figure \ref{fig08}. The comparison of gravity datasets gives additional information on the distribution of major geological and seafloor structures with regard to the variations of geophysical fields. 

\begin{figure}[H]
\centering
	\hspace{100pt}\includegraphics[width=14.0 cm]{Fig_08.jpg}
\caption{Vertical gravity gradient over Kerguelen Plateau. Mapping: GMT. Map source: author. \label{fig08}}
\end{figure}

Moreover, the comparison of the vertical gradient and free-air gravity map (Figure \ref{fig07} and Figure \ref{fig08}) with the topographic map (Figure \ref{fig01}) illustrates the effects of seafloor structures and their distribution of the oceanic bed on gravity and showed that the highest values correspond to the Kerguelen Plateau and other rises, while lower values are generally associated with depressions.

%----------- subsection ----------->
\subsection{Magnetic anomalies over Kerguelen Plateau}

The anomaly of the magnetic intensity at an altitude of 5 km above mean sea level over Kerguelen Plateau is shown in Figure \ref{fig09}. High heterogeneity in the geophysical data is related to the past volcanism over the Kerguelen Plateau which included the voluminous basaltic flooding originated from a deep hot spot as an asthenospheric source of mantle plume product as a result of slab dynamics and tectonic plate movements in the southwestern segment of the Indian Ocean. In this regard, combining the data from the WDMAM and EMAG2 (Figure \ref{fig09} and Figure \ref{fig10}) data on terrestrial gravity fields (Figure \ref{fig07} and Figure \ref{fig08}) for comparison with maps on oceanic crust development (Figures \ref{fig02} to \ref{fig05}) presents an integrated geophysical analysis. 

\begin{figure}[H]
\centering
	\hspace{100pt}\includegraphics[width=14.0 cm]{Fig_09.jpg}
\caption{Patterns of the marine and terrestrial airborne magnetic anomaly over the Kerguelen Plateau on a 3-min resolution grid of WDMAM, south-west Indian Ocean. Map source: author. \label{fig09}}
\end{figure}

Crustal volume contributes to the decreased amplitude of the magnetic anomaly around the Kerguelen, as can be seen as dark blue areas in Figure \ref{fig10} with lower values of around -500 mGal. The analysis of the magnetic anomaly patterns in the SW Indian Ocean supports the hypothesis of the spreading seafloor with variations in the oceanic crustal blocks movements, as reported earlier \cite{LePichon}, which is visible on different magnetic pattern over the mid-ocean ridge. The comparison of Figure \ref{fig10} with Figure \ref{fig05} reveals that the age uncertainties for grid cells coincide with the marine magnetic anomaly identified around the Kerguelen, which is also observed for the conjugate ridge flanks (Figure \ref{fig03} and Figure \ref{fig10}).

In fact, the Earth Magnetic Anomaly Grid (EMAG2) enables to evaluate the relationship in magnetic-gravity field aspects as descriptors besides the to the traditional analysis of gravity and magnetic anomalies. Since the magnetic anomalies result from geologic and topographic features that increase or depress the local magnetic field, understanding the correlation in geophysical phenomena with topography is essential. Thus, the analysis of local patterns of magnetic anomalies in the southwest Indian Ocean shows patterns of magnetic anomalies related to the formation of oceanic crust and seafloor spreading, as well as the and existing subduction zones of oceanic crust. Moreover, the age of the of oceanic crust and the spreading rates by accretion of land areas and large-scale volcanism of the Kerguelen Plateau are associated with past geological development of the region and spreading of the seafloor of the Indian Ocean, as noted earlier \cite{Maus2009}.

The EMAG-2 and WDMAM grids used for plotting the magnetic anomalies show various levels of details of the grids which enable to reveal the subsurface structure of the seafloor around the Kerguelen Plateau and the composition of the Earth's crust in the southwest Indian Ocean. Here, the magnetic fabric data correlate with the hotspot activities and active volcanism which is notable over the central and northern parts of the archipelago \cite{HENRY2003153,Xu,Pettersen}. Moreover, values of the intermediate crustal thickness in the oceanic crust beneath the Kerguelen Plateau and large-amplitude magnetic anomalies over the archipelago evidence that the plateau is oceanic in origin and is associated with the plate volcanism caused by the tectonic plate activity, which is also reported in earlier studies \cite{Direen,Charvis1995}. Therefore, the distinct magnetic patterns visible in maps in Figure \ref{fig10} corresponds to the high hotspot activity and associated ava flows. 

\begin{figure}[H]
\centering
	\hspace{100pt}\includegraphics[width=14.0 cm]{Fig_10.jpg}
\caption{Marine and Earth airborne magnetic anomaly grid based on EMAG-2 over the Kerguelen Plateau region. Black areas signify "no data" in the original grid. Map source: author. \label{fig10}}
\end{figure}

The effects from deeper mass, asthenospheric upwelling and convection from mantle plumes as geodynamic processes and affect magnetic properties of the surface. Moreover, the analysis of data enables to better understand the impact of geophysical setting on  on the distribution of positive (over the Kerguelen Islands) and negative undulation (eastern regions of the study area and south-west Australia), which enables to evaluate the variations of geophysical grids using the comparative analysis of maps. In this way, cartographic representation of the geophysical and magnetic datasets provides advanced methods to obtain data on the seafloor formation and interrelated geophysical process.

%----------- subsection ----------->
\subsection{Geoid models}
The geopotential model over the Kerguelen based on the EGM96 dataset is shown in Figure \ref{fig11}. The variations of geoid over the Kerguelen Plateau point at the existing isostatic compensation of the archipelago due to the low-density mantle. Thus, high anomalies in the geoid level above the Kerguelen Plateau are explained by a high volcanism associated to the ridge formation on the hot lithosphere, which is reflected in the extraordinary thickness below the Kerguelen with geoid values of above 40 m (red areas in Figure \ref{fig11}) which exceeds normal thickness of the oceanic crust. This proves previous investigations on geoid in the Kerguelen Archipelago that reported anomalous thickness in this area \cite{RECQ1987301}. Such isostatic compensation is associated with the anomalously high geoid values and is related to the topographically rough elevated surfaces in regions of active tectonic uplift which identifies processes in upper mantle \cite{Coblentz}. 

\begin{figure}[H]
\centering
	\hspace{100pt}\includegraphics[width=14.0 cm]{Fig_11.jpg}
\caption{Geoid model based on EGM-96 over the Kerguelen Plateau region, east Antarctic and south-west Indian Ocean. Mapping: GMT. Map source: author. \label{fig11}}
\end{figure}

A comparison of the geoid map with topography (sf. Figures \ref{fig01} and \ref{fig11}) implies correlation between the continued geodynamic processes in the south-west Indian Ocean and the topographic structure of the Kerguelen Archipelago. Moreover, this proves a high positive correlation between the geoid height and deep structure of the seafloor topography, which is noted earlier for the regions of large plateaus and swells \cite{SandwellMacKenzie}. 

%----------- section ----------->
\section{Conclusions}

Global surveillance of seafloor using altimeter satellites and gravity enables to reveal geophysical anomalies. Furthermore, integrated data on magnetic field intensity, bathymetry and geodynamics of the deep seafloor enables evaluation of the associations between these processes. Global surveillance of seafloor delivers ever-increasing volumes of high resolution datasets. However, the scientific application of these data in ever-higher spatial resolution and precision requires advanced tools  for their effective automated analysis. GMT scripts toolset has proven to be effective for seafloor bathymetric and geophysical mapping. Diverse seabed features in different scales and resolution can be visualised and mapped for detection of coherency between the magnetic anomalies and geophysical patterns and their relationship with ocean crust geodynamics reflected in present bathymetry. Visualisation of such data provides an extensive near-global coverage of seafloor features. 

The direct surveys of seafloor observation are expensive to generate bathymetric data and require the use of complex costly equipment such as multi-beam echo-sounding systems for data sampling, generating and storage. At the same time, the reality calls for efficient datasets for visualisation and analysis. Using high-resolution geophysical datasets provided by NOAA and USGS processed by GMT advanced scripting approach develops a cartographic processing pipeline for fast, automated and accurate mapping of seafloor. A balance between the topographic gradients and geophysical grids well illustrates the links between the seafloor patterns and the structure of the oceanic crust and the processes in the lithospheric mantle. From the cartographic point of view, this proves that the effective analysis of the geologic-geophysical datasets in the Earth sciences is not confined to using selected data (e.g., topographic maps) as a single parameter but also to properly select several datasets since the processes of deep mantle reflected in the geophysical data define the process of seafloor formation. 

The cartographic approach demonstrated in this study makes possible data assimilation and extend it not only on the Indian Ocean but also towards other regions of the World Ocean. For instance, the Pacific Ocean is notable for a rich seafloor texture features such as vast abyssal plains, mid-ocean ridges, oceanic trenches, numerous seamounts and continental shelf which provide a reliable source of information for matching potential correlation between the geophysical and magnetic anomalies and bathymetric responses as heterogeneous seafloor pattern. With this regard, the use of the GMT for seafloor mapping in other regions of the World Ocean makes an ideal case for validating our cartographic scripting approach presented in the above methodological scripting framework. 

The above series of maps and the comparison of the geophysical setting over the Kerguelen Plateau performed using the GMT conclude that scripts techniques outperform the GIS approach in terms of automation of the cartographic workflow. Besides, the compactness of the code using the GMT syntax which can be reapplied with modified setting. When handling the cartographic elements on maps and adjusting the parameters of projection, the main drawback of the GMT is that it needs to tune these parameters in advance while writing the code. This is because it cannot visualise the map prior to the executed script being a completely console-based program. Furthermore, the GMT cannot discard the redundant features once plotted on a map. Therefore, correcting the map requires running the script anew. In GIS approach, for example, user's setting may change the appearance of map layout so that new colour palettes and details of the bathymetric plotting with regard to the isoline density and consistency may become redundant and be corrected on the fly. In contrast, as a console-based application, GMT requires modification of these parameters directly in the code lines of the script to handle cartographic issues. 

%----------- section ----------->

\funding{The publication was funded by the Editorial Office of \emph{Geomatics}, Multidisciplinary Digital Publishing Institute (MDPI), by providing 100\% discount for the APC of this manuscript.}

\institutionalreview{Not applicable.}

\informedconsent{Not applicable.}

\dataavailability{Not applicable} 

\acknowledgments{The author sincerely thanks the two anonymous reviewers for critical review of this paper,  constructive comments and useful suggestions which improved the initial version of this manuscript. The author is grateful to the Editorial Office of the Geomatics MDPI for support of this publication.}

\conflictsofinterest{The authors declare no conflict of interest.} 

\abbreviations{Abbreviations}{
The following abbreviations are used in this manuscript:\\

\noindent 
\begin{tabular}{@{}ll}
DEM & Digital Elevation Model \\
EMAG-2 & Earth Magnetic Anomaly Grid 2 \\
\hl{EML} & \hl{ERDAS Macro Language} \\
\hl{ESRI} & \hl{Environmental Systems Research Institute} \\
GDAL & Geospatial Data Abstraction Library \\
GRACE & Gravity Recovery and Climate Experiment \\
GMT & Generic Mapping Tools \\
GRASS GIS & Geographic Resources Analysis Support System \\
GUI & Graphical User Interface \\
NOAA & National Oceanic and Atmospheric Administration \\
SRTM & Shuttle Radar Topography Mission \\
USGS & United States Geological Survey \\
UTM & Universal Transverse Mercator \\
WGM & World Gravity Model \\
WDMAM & World Digital Magnetic Anomaly Map\\
\end{tabular}
}

\appendixtitles{yes} % Leave argument "no" if all appendix headings stay EMPTY (then no dot is printed after "Appendix A"). If the appendix sections contain a heading then change the argument to "yes".
\appendixstart
\appendix
\section[\appendixname~\thesection]{GMT scripts}\label{App}

\subsection[\appendixname~\thesubsection]{GMT script for mapping topography of the Kerguelen region in Lambert Azimuthal Equal-Area projection}

\begin{lstlisting}[language=bash,caption=Mapping topography of the Kerguelen region in Lambert Azimuthal Equal-Area projection,style=mystyle,label={lst01}]
# Subset study area
grdcut ETOPO1_Ice_g_gmt4.grd -R-40/150/-70/-20 -Gkgl_relief.nc
gdalinfo ss_relief.nc -stats
# Minimum=-8239.000, Maximum=6392.000
# Color palette
gmt makecpt -Cgray.cpt -V -T-8239/6392 > myocean.cpt
# Generate a file
ps=Bathymetry_Kgl.ps
gmt grdimage kgl_relief.nc -Cmyocean.cpt -R-25/-65/101/-10r -JA55/-50/7.5i -P -I+a15+ne0.75 -Xc -K > $ps
# Add shorelines
gmt grdcontour kgl_relief.nc -R -J -C1000 -Wthinnest,blue -O -K >> $ps
# Add grid
gmt psbasemap -R -J \
    -Bpxg10f5a5 -Bpyg10f5a5 -Bsxg5 -Bsyg5 \
    -B+t"Topographic map of the Kerguelen Plateau" \
    -Lx15.0c/-1.3c+c318/-57+w2000k+l"Scale (km) at 60\232E 50\232S"+f \
    -UBL/-5p/-40p -O -K >> $ps
# Texts
gmt pstext -R -J -X0.0c -Y0.0c -N -O -K \
    -F+jTL+f9p,Helvetica,white+jLB >> $ps << EOF
66.0 -62.0 Enderby
66.5 -63.0 Basin
EOF
# Other text are added likewise
# Add legend
gmt psscale -Dg-27.0/-60+w15.4c/0.4c+v+ml -R -J -Cmyocean.cpt \
    -Bg1000f200a2000+l"Color scale: geo global bathymetry/topography relief [R=-8239/6392, H=0, C=RGB]" \
    -I0.2 -By+lm -O -K >> $ps
# Add GMT logo
gmt logo -Dx5.5/-2.2+o0.1i/0.1i+w2c -O -K >> $ps
# Add subtitle
gmt pstext -R0/10/0/15 -JX10/10 -X0.5c -Y13.0c -N -O \
    -F+f10p,Palatino-Roman,black+jLB >> $ps << EOF
4.0 6.1 ETOPO1 global terrain model, 1 arc min resolution grid
2.0 5.5 Lambert Azimuthal Equal-Area projection. Central meridian 55\232E, standard parallel 50\232S
EOF
# Convert to image file using GhostScript
gmt psconvert Bathymetry_Kgl.ps -A1.0c -E720 -Tj -Z
\end{lstlisting}

\subsection[\appendixname~\thesubsection]{GMT script for mapping age, spreading rates, spreading asymmetry and age uncertainty of the ocean crust}

\begin{lstlisting}[language=bash,caption=Mapping age of the oceanic lithosphere over Kerguelen Plateau and SW Indian Ocean,style=mystyle,label={lst02}]
gmt grdcut age.3.2.nc -R-40/150/-70/-10 -Gker_age.tif
gdalinfo ker_age.tif -stats
# Minimum=0.000, Maximum=16001.000, Mean=5895.761, StdDev=4042.817
# Color palette
# gmt makecpt -Cjet -T14/16000.000 > age.cpt
gmt makecpt -Cwysiwyg -T14/16000.000 > age.cpt
#gmt makecpt --help
# Generate a file
ps=Ker_age.ps
gmt grdimage ker_age.tif -Cage.cpt -R-25/-65/101/-10r -JA55/-50/7.5i -P -I+a15+ne0.75 -Xc -K > $ps
# add grid
gmt psbasemap -R -J \
    -Bpxg10f5a10 -Bpyg10f5a10 -Bsxg5 -Bsyg5 \
    -B+t"Age of Oceanic Lithosphere: Kerguelen Plateau and SW Indian Ocean" -O -K >> $ps
# Legend
gmt psscale -Dg-30/-58+w15.4c/0.4c+v+ml+e -R -J -Cage.cpt \
    -Bg1000f200a1000+l"Color scale: 'wysiwyg' [R=0/6000, H=0, C=RGB]" \
    -I0.2 -By+l"M.Y." -O -K >> $ps   
# Scale, directional rose
gmt psbasemap -R -J \
    -Tdx0.8c/10.3c+w0.3i+f2+l+o0.15i \
    -Lx16.0c/-1.6c+c318/-57+w2000k+l"Scale (km) at 55\232E 50\232S"+f \
    -UBL/-5p/-40p -O -K >> $ps
# GMT logo
gmt logo -Dx5.5/-2.2+o0.1i/0.1i+w2c -O -K >> $ps
# Subtitle
gmt pstext -R0/10/0/15 -JX10/10 -X0.5c -Y13.0c -N -O \
    -F+f12p,0,black+jLB >> $ps << EOF
0.0 5.9 Lambert Azimuthal Equal-Area projection. Central meridian 55\232E, standard parallel 50\232S
0.0. 6.6 Age of the ocean crust, 2 arc min netCDF grid v.3., according to adjacent seafloor isochrons
EOF
# Convert to image file using GhostScript
gmt psconvert Ker_age.ps -A1.7c -E720 -Tj -Z
\end{lstlisting}

\subsection[\appendixname~\thesubsection]{GMT script for mapping asymmetries in crustal accretion on conjugate ridge flanks: Kerguelen Plateau and SW Indian Ocean}

\begin{lstlisting}[language=bash,caption=Mapping asymmetries in crustal accretion on conjugate ridge flanks: Kerguelen Plateau and SW Indian Ocean,style=mystyle,label={lst03}]
exec bash
# Cut off raster image
gmt grdcut asym.3.2.nc -R-40/150/-70/-10 -Gker_asym.tif
gdalinfo ker_asym.tif -stats
# Minimum=14.000, Maximum=10000.000, Mean=5491.073, StdDev=1628.733
# Make color palette
gmt makecpt -Cturbo -T14/10000.000 > asym.cpt
# Generate a file
ps=Ker_asym.ps
gmt grdimage ker_asym.tif -Casym.cpt -R-25/-65/101/-10r -JA55/-50/7.5i -P -I+a15+ne0.75 -Xc -K > $ps
# Grid
gmt psbasemap -R -J \
    -Bpxg10f5a10 -Bpyg10f5a10 -Bsxg5 -Bsyg5 \
    -B+t"Asymmetries in crustal accretion (%) on conjugate ridge flanks: Kerguelen Plateau and SW Indian Ocean" -O -K >> $ps
# Legend
gmt psscale -Dg-30/-58+w15.4c/0.4c+v+ml+e -R -J -Casym.cpt \
    -Bg1000f200a1000+l"Color scale: 'no_green' [R=0/6000, H=0, C=RGB]" \
    -I0.2 -By+lm -O -K >> $ps
# Scale, directional rose
gmt psbasemap -R -J \
    -Tdx0.8c/10.3c+w0.3i+f2+l+o0.15i \
    -Lx16.0c/-1.6c+c318/-57+w2000k+l"Scale (km) at 55\232E 50\232S"+f \
    -UBL/-5p/-40p -O -K >> $ps
# GMT logo
gmt logo -Dx5.5/-2.2+o0.1i/0.1i+w2c -O -K >> $ps
# Subtitle
gmt pstext -R0/10/0/15 -JX10/10 -X0.5c -Y13.0c -N -O \
    -F+f12p,0,black+jLB >> $ps << EOF
0.0 5.9 Lambert Azimuthal Equal-Area projection. Central meridian 55\232E, standard parallel 50\232S
0.0. 6.6 Age, spreading rates and spreading asymmetry of the ocean crust, 2 arc min netCDF grid v.3
EOF
# Convert to image file using GhostScript
gmt psconvert Ker_asym.ps -A1.7c -E720 -Tj -Z
\end{lstlisting}

\subsection[\appendixname~\thesubsection]{GMT script for mapping spreading half rates of the oceanic lithosphere in Kerguelen Plateau and SW Indian Ocean}

\begin{lstlisting}[language=bash,caption=Mapping spreading half rates of the oceanic lithosphere in Kerguelen Plateau and SW Indian Ocean,style=mystyle,label={lst04}]
# GMT set up
gmt set FORMAT_GEO_MAP=dddF \
    MAP_FRAME_PEN=dimgray \
    MAP_FRAME_WIDTH=0.1c \
    MAP_TITLE_OFFSET=0.5c \
    MAP_ANNOT_OFFSET=0.1c \
    MAP_TICK_PEN_PRIMARY=thinner,dimgray \
    MAP_GRID_PEN_PRIMARY=thin,white \
    MAP_GRID_PEN_SECONDARY=thinnest,white \
    FONT_TITLE=12p,0,black \
    FONT_ANNOT_PRIMARY=7p,0,dimgray \
    FONT_LABEL=7p,0,dimgray
#Overwrite defaults of GMT
gmtdefaults -D > .gmtdefaults
exec bash
# Cut off raster image
gmt grdcut rate.3.6.nc -R-40/150/-70/-10 -Gker_rate.tif
gdalinfo ker_rate.tif -stats
# Minimum=0.000, Maximum=15000.000, Mean=3175.093, StdDev=2125.065
# Make color palette
gmt makecpt -Cf-30-31-32.cpt -T0/15000/250 -N -Iz > age.cpt
# Generate a file
ps=Ker_rate.ps
gmt grdimage ker_rate.tif -Cage.cpt -R-25/-65/101/-10r -JA55/-50/7.5i -P -I+a15+ne0.75 -Xc -K > $ps
# Add grid
gmt psbasemap -R -J \
    -Bpxg10f5a10 -Bpyg10f5a10 -Bsxg5 -Bsyg5 \
    -B+t"Spreading Half Rates of the Oceanic Lithosphere in Kerguelen Plateau and SW Indian Ocean" -O -K >> $ps
gmt psxy -R -J ridge.gmt -Sf0.5c/0.15c+l+t -Wthin,yellow -Gpurple -O -K >> $ps
gmt psxy -R -J TP_Indian.txt -L -Wthickest,red -O -K >> $ps
gmt psxy -R -J TP_Australian.txt -L -Wthickest,red -O -K >> $ps
gmt psxy -R -J TP_African.txt -L -Wthickest,red -O -K >> $ps
# tectonic slab contours
gmt psxy -R -J GSFML_SF_FZ_KM.gmt -Wthick,cyan2 -O -K >> $ps
gmt psxy -R -J GSFML_SF_FZ_RM.gmt -Wthick,cyan2 -O -K >> $ps
# transform faults
gmt psxy -R -J transform.gmt -Sc0.05c -Ggreen -Wthick,deeppink1 -O -K >> $ps 
# Add legend
gmt psscale -Dg-30/-58+w15.4c/0.4c+v+ml+e -R -J -Cage.cpt \
    -Bg1000f200a1000+l"Color scale: 'f-30-31-32.cpt' from Gnuplot [R=0/15000, H=0, C=RGB]" \
    -I0.2 -By+l"mm/yr." -O -K >> $ps 
# Add scale, directional rose
gmt psbasemap -R -J \
    -Tdx0.8c/10.3c+w0.3i+f2+l+o0.15i \
    -Lx16.0c/-1.6c+c318/-57+w2000k+l"Scale (km) at 55\232E 50\232S"+f \
    -UBL/-5p/-40p -O -K >> $ps
# Add GMT logo
gmt logo -Dx5.5/-2.2+o0.1i/0.1i+w2c -O -K >> $ps
# Add subtitle
gmt pstext -R0/10/0/15 -JX10/10 -X0.5c -Y13.0c -N -O \
    -F+f12p,0,black+jLB >> $ps << EOF
0.0 5.9 Lambert Azimuthal Equal-Area projection. Central meridian 55\232E, standard parallel 50\232S
0.0. 6.6 Global Seafloor Fabric and Magnetic Lineation Data (GSFML) are shown by cyan lines
EOF
# Convert to image file using GhostScript
gmt psconvert Ker_rate.ps -A1.7c -E720 -Tj -Z
\end{lstlisting}

\subsection[\appendixname~\thesubsection]{GMT script for mapping crustal age uncertainty of the oceanic lithosphere in Kerguelen Plateau and SW Indian Ocean}

\begin{lstlisting}[language=bash,caption=Mapping crustal age uncertainty of the oceanic lithosphere in Kerguelen Plateau and SW Indian Ocean,style=mystyle,label={lst05}]
#!/bin/sh
# Age, spreading rates and spreading asymmetry of the ocean crust, 2 arc min netCDF grid v 3
# GMT modules: gmtset, gmtdefaults, grdcut, makecpt, grdimage, psscale, grdcontour, psbasemap, gmtlogo, psconvert
exec bash
# Cut off raster image
gmt grdcut ageerror.3.2.nc -R-40/150/-70/-10 -Gker_error.tif
gdalinfo ker_error.tif -stats
# Minimum=8.000, Maximum=1313.000, Mean=196.616, StdDev=156.744
# Color palette
gmt makecpt -Ccyan-magenta-yellow-white.cpt -T8/800.000 -N > age.cpt
# Generate a file
ps=Ker_error.ps
gmt grdimage ker_error.tif -Cage.cpt -R-25/-65/101/-10r -JA55/-50/7.5i -P -I+a15+ne0.75 -Xc -K > $ps
# Add grid
gmt psbasemap -R -J -Bpxg10f5a10 -Bpyg10f5a10 -Bsxg5 -Bsyg5 \
    -B+t"Crustal Age Uncertainty of the Oceanic Lithosphere in Kerguelen Plateau and SW Indian Ocean" -O -K >> $ps
# gmt psxy -R -J ridge.gmt -Sf0.5c/0.15c+l+t -Wthin,yellow -Gpurple -O -K >> $ps
gmt psxy -R -J TP_Indian.txt -L -Wthickest,red -O -K >> $ps
gmt psxy -R -J TP_Australian.txt -L -Wthickest,red -O -K >> $ps
gmt psxy -R -J TP_African.txt -L -Wthickest,red -O -K >> $ps
# tectonic slab contours
# gmt psxy -R -J GSFML_SF_FZ_KM.gmt -Wthick,slateblue3 -O -K >> $ps
gmt psxy -R -J GSFML_SF_FZ_KM.gmt -Wthick,darkmagenta -O -K >> $ps
gmt psxy -R -J GSFML_SF_FZ_RM.gmt -Wthick,darkmagenta -O -K >> $ps
# transform faults
gmt psxy -R -J transform.gmt -Sc0.05c -Ggreen -Wthick,deeppink1 -O -K >> $ps  
# Legend
gmt psscale -Dg-30/-58+w15.4c/0.4c+v+ml+e -R -J -Cage.cpt \
    --FONT_LABEL=10p,0,dimgray --FONT_ANNOT_PRIMARY=10p,0,black \
    -Bg100f20a100+l"Color scale: 'jet' [R=0/15000, H=0, C=RGB]" \
    -I0.2 -By+l"m/yr." -O -K >> $ps
# Scale, directional rose
gmt psbasemap -R -J \
    --FONT=10p,0,black --MAP_TITLE_OFFSET=0.3c \
    -Tdx0.8c/10.3c+w0.3i+f2+l+o0.15i \
    -Lx16.0c/-1.6c+c318/-57+w2000k+l"Scale (km) at 55\232E 50\232S"+f \
    -UBL/-5p/-40p -O -K >> $ps
# GMT logo
gmt logo -Dx5.5/-2.2+o0.1i/0.1i+w2c -O -K >> $ps
# Add subtitle
gmt pstext -R0/10/0/15 -JX10/10 -X0.5c -Y13.0c -N -O \
    -F+f12p,0,black+jLB >> $ps << EOF
0.0 5.9 Lambert Azimuthal Equal-Area projection. Central meridian 55\232E, standard parallel 50\232S
0.0. 6.6 Global Seafloor Fabric and Magnetic Lineation Data (GSFML) are shown by cyan lines
EOF
# Convert to image file using GhostScript
gmt psconvert Ker_error.ps -A1.7c -E720 -Tj -Z
\end{lstlisting}

\subsection[\appendixname~\thesubsection]{GMT script for sediment thickness over the Kerguelen Plateau and SW Indian Ocean}

\begin{lstlisting}[language=bash,caption=Mapping sediment thickness over the Kerguelen Plateau and SW Indian Ocean,style=mystyle,label={lst06}]
exec bash
#gmt grdcut GlobSed-v2.nc -R-25/-65/101/-10r -Gker_sed.nc
gmt grdcut GlobSed-v2.nc -R-40/150/-70/-10 -Gker_sed.nc
gdalinfo ker_sed.nc -stats
# Minimum=0.000, Maximum=8116.000, Mean=530.333, StdDev=728.622
# Make color palette
gmt makecpt -Cno_green.cpt -V -T0/6000 > myocean.cpt
# gmt makecpt --help
# Generate a file
ps=SedThick_Kgl.ps
gmt grdimage ker_sed.nc -Cmyocean.cpt -R-25/-65/101/-10r -JA55/-50/7.5i -P -I+a15+ne0.75 -Xc -K > $ps
# Add shorelines
gmt grdcontour ker_sed.nc -R -J -C200 -A200+f10p,25,black -Wthinner,dimgray -O -K >> $ps
# Add grid
gmt psbasemap -R -J \
    -Bpxg10f5a10 -Bpyg10f5a10 -Bsxg5 -Bsyg5 \
    -B+t"Sediment thickness over East Antarctic, Kerguelen Plateau and SW Indian Ocean" \
    -Lx15.0c/-1.5c+c318/-57+w2000k+l"Scale (km) at 60\232E 50\232S"+f \
    -UBL/-5p/-40p -O -K >> $ps
# Texts
# Add legend
gmt psscale -Dg-33/-59+w15.4c/0.4c+v+ml+e -R -J -Cmyocean.cpt \
    -Bg1000f50a1000+l"Color scale: 'no_green' [R=0/6000, H=0, C=RGB]" \
    -I0.2 -By+lm -O -K >> $ps
# Add GMT logo
gmt logo -Dx5.5/-2.2+o0.1i/0.1i+w2c -O -K >> $ps
# Add subtitle
gmt pstext -R0/10/0/15 -JX10/10 -X0.5c -Y13.0c -N -O \
    -F+f12p,0,black+jLB >> $ps << EOF
3.0 6.1 GlobSed: Total Sediment Thickness Version 3, 5 arc minute grid
EOF
# Convert to image file using GhostScript
gmt psconvert SedThick_Kgl.ps -A1.5c -E720 -Tj -Z
\end{lstlisting}

\subsection[\appendixname~\thesubsection]{GMT scripts for geophysical mapping over the Kerguelen Plateau}

\begin{lstlisting}[language=bash,caption=Mapping free-air gravity anomalies over Kerguelen Plateau and SW Indian Ocean,style=mystyle,label={lst07}]
exec bash
gmt img2grd grav_27.1.img -R40/110/-70/-20 -Ggrav_Ker.grd -T1 -I1 -E -S0.1 -V
gdalinfo grav_Ker.grd -stats
gmt makecpt -Chaxby -V -T-100/150 > myocean.cpt
ps=Gravity_Kgl.ps
gmt grdimage grav_Ker.grd -Cmyocean.cpt -R40/110/-70/-20 -JU43/6.0i -P -I+a15+ne0.75 -Xc -K > $ps
gmt grdcontour grav_Ker.grd -R -J -C50 -A50 -Wthinner -O -K >> $ps
gmt psbasemap -R -J \
    -Bpxg10f5a10 -Bpyg10f5a5 -Bsxg5 -Bsyg5 \
    -B+t"Free-air gravity anomaly on Kerguelen Plateau" \
    -Lx7.5c/-1.3c+c318/-57+w2000k+l"UTM projection, Zone 43. Scale (km)"+f \
    -UBL/10p/-40p -O -K >> $ps
gmt psscale -Dg41/-12.5+w15.4c/0.4c+ml+h+e -R -J -Cmyocean.cpt \
    -Bg20f2a20+l"Color scale: haxby [R=-100/547, H=0, C=RGB]" \
    -I0.2 -By+lm -O -K >> $ps
gmt logo -Dx7.5/-2.2+o0.1i/0.1i+w2c -O -K >> $ps
gmt pstext -R0/10/0/15 -JX10/10 -X0.5c -Y13.0c -N -O \
    -F+f10p,0,black+jLB >> $ps << EOF
1.0 1.3 Global gravity grid from CryoSat-2 and Jason-1, 1 min resolution, SIO, NOAA, NGA
EOF
gmt psconvert Gravity_Kgl.ps -A1.0c -E720 -Tj -Z
\end{lstlisting}

\begin{lstlisting}[language=bash,caption=Mapping vertical gravity gradient over Kerguelen Plateau and SW Indian Ocean,style=mystyle,label={lst08}]
exec bash
gmt img2grd curv_27.1.img -R40/110/-70/-20 -Ggravvert_Ker.grd -T1 -I1 -E -S0.1 -V
gdalinfo gravvert_Ker.grd -stats
gmt makecpt -Cmag.cpt -T-40/40 > colors.cpt
ps=Gravity_Kgl_vert.ps
gmt grdimage gravvert_Ker.grd -Ccolors.cpt -R40/110/-70/-20 -JU43/6.0i -P -I+a15+ne0.75 -Xc -K > $ps
gmt grdcontour gravvert_Ker.grd -R -J -C100 -A100 -Wthinner -O -K >> $ps
gmt psbasemap -R -J \
    -Bpxg10f5a10 -Bpyg10f5a5 -Bsxg5 -Bsyg5 \
    -B+t"Vertical gravity gradient over Kerguelen Plateau" \
    -Lx7.5c/-1.3c+c318/-57+w2000k+l"UTM projection, Zone 43. Scale (km)"+f \
    -UBL/10p/-40p -O -K >> $ps
gmt psscale -Dg41/-12.5+w15.4c/0.4c+ml+h+e -R -J -Ccolors.cpt \
    -Bg20f2a20+l"Color scale: mag [R=-40/40, H=0, C=RGB]" \
    -I0.2 -By+l"mGal/m" -O -K >> $ps
gmt logo -Dx7.5/-2.2+o0.1i/0.1i+w2c -O -K >> $ps
# Add subtitle
gmt pstext -R0/10/0/15 -JX10/10 -X0.5c -Y13.0c -N -O \
    -F+f10p,0,black+jLB >> $ps << EOF
1.0 1.3 Global gravity grid derived from satellite altimetry (CryoSat-2 and Jason-1: NOAA, NGA)
EOF
gmt psconvert Gravity_Kgl_vert.ps -A1.0c -E720 -Tj -Z
\end{lstlisting}

\subsection[\appendixname~\thesubsection]{GMT scripts for mapping magnetic anomalies over the Kerguelen Plateau}

\begin{lstlisting}[language=bash,caption=Mapping magnetic anomalies by WDMAM over Kerguelen Plateau,style=mystyle,label={lst09}]
exec bash
gmt grdcut @earth_wdmam_03m -R50/100/-70/-30 -Gker_mag.nc
gdalinfo ker_mag.nc -stats
gmt makecpt -Cmag.cpt -V -T-500/500 > myocean.cpt
ps=Magnet_Kgl.ps
gmt grdimage ker_mag.nc -Cmyocean.cpt -R50/100/-70/-30 -JM6.5i -P -I+a15+ne0.75 -Xc -K > $ps
gmt grdcontour ker_mag.nc -R -J -C200 -A1+f10p,25,black -Wthinner,dimgray -O -K >> $ps
gmt psbasemap -R -J \
    -Bpxg10f5a5 -Bpyg10f5a5 -Bsxg5 -Bsyg5 \
    -B+t"Marine and Earth airborne Magnetic Anomaly Grid on Kerguelen Plateau" \
    -Lx14.0c/-3.0c+c318/-57+w1000k+l"Mercator projection. Scale (km)"+f \
    -UBL/-5p/-80p -O -K >> $ps
gmt psscale -Dg50/-71.5+w16.0c/0.4c+h+ml+e -R -J -Cmyocean.cpt \
    -Bg100f10a50+l"Color scale: mag (Colors for magnetic anomaly maps), [C=RGB]" \
    -I0.2 -By+l"nT" -O -K >> $ps
gmt logo -Dx6.5/-3.8+o0.1i/0.1i+w2c -O -K >> $ps
gmt pstext -R0/10/0/15 -JX10/10 -X0.5c -Y13.0c -N -O \
    -F+f11p,0,black+jLB >> $ps << EOF
1.2 15.8 WDMAM (World Digital Magnetic Anomaly Map), 3 arc min resolution grid
EOF
gmt psconvert Magnet_Kgl.ps -A1.0c -E720 -Tj -Z
\end{lstlisting}

\begin{lstlisting}[language=bash,caption=Mapping magnetic anomalies by EMAG2 over Kerguelen Plateau,style=mystyle,label={lst10}]
exec bash
gmt grdcut EMAG2_V2.grd -R50/100/-70/-30 -Gker_mag.nc
gdalinfo ker_mag.nc -stats
gmt makecpt -Cmag.cpt -V -T-500/500 > myocean.cpt
ps=Magnet_Kgl.ps
gmt grdimage ker_mag.nc -Cmyocean.cpt -R50/100/-70/-30 -JM6.5i -P -I+a15+ne0.75 -Xc -K > $ps
gmt grdcontour ker_mag.nc -R -J -C200 -A1+f10p,25,black -Wthinner,dimgray -O -K >> $ps
gmt psbasemap -R -J \
    -Bpxg10f5a5 -Bpyg10f5a5 -Bsxg5 -Bsyg5 \
    -B+t"Detailed marine and Earth EMAG-2 airborne Magnetic Anomaly Grid on Kerguelen Plateau" \
    -Lx14.0c/-3.0c+c318/-57+w1000k+l"Mercator projection. Scale (km)"+f \
    -UBL/-5p/-80p -O -K >> $ps
gmt psscale -Dg50/-71.5+w16.0c/0.4c+h+ml+e -R -J -Cmyocean.cpt \
    -Bg100f10a50+l"Color scale: mag (Colors for magnetic anomaly maps), [C=RGB]" \
    -I0.2 -By+l"nT" -O -K >> $ps
gmt logo -Dx6.5/-3.8+o0.1i/0.1i+w2c -O -K >> $ps
gmt pstext -R0/10/0/15 -JX10/10 -X0.5c -Y13.0c -N -O \
    -F+f11p,0,black+jLB >> $ps << EOF
1.2 15.8 EMAG-2 (Earth Magnetic Anomaly Grid, 2 arc min resolution)
EOF
gmt psconvert Magnet_Kgl.ps -A1.0c -E720 -Tj -Z
\end{lstlisting}

\subsection[\appendixname~\thesubsection]{GMT scripts for mapping geoid anomalies over the Kerguelen Plateau}

\begin{lstlisting}[language=bash,caption=Mapping geoid anomalies in Kerguelen Plateau by EGM96,style=mystyle,label={lst11}]
exec bash
gdalinfo geoid.egm96.grd -stats
gmt grd2cpt geoid.egm96.grd -Cjet > geoid.cpt
ps=Geoid_Ker.ps
gmt grdimage geoid.egm96.grd -I+a45+nt1 -R40/110/-70/-20 -JQ7.5i -Cgeoid.cpt -P -K > $ps
gmt grdcontour geoid.egm96.grd -R -J -C2 -A4+f10p,25,black -Wthinner,dimgray -O -K >> $ps
gmt psbasemap -R -J \
    -Bpxg10f5a10 -Bpyg10f5a10 -Bsxg5 -Bsyg5 \
    -B+t"Geoid gravitational model EGM96 over East Antarctic, Kerguelen Plateau and SW Indian Ocean" \
    -Lx16.0c/-1.5c+c318/-57+w1000k+l"Cylindrical equidistant projection. Scale (km)"+f \
    -UBL/-5p/-40p -O -K >> $ps
gmt psscale -Dg33/-70+w13.4c/0.4c+v+ml+e -R -J -Cgeoid.cpt \
    -Bg10f1a10+l"Color scale: jet [C=RGB]" \
    -I0.2 -By+lm -O -K >> $ps
gmt logo -Dx7.0/-2.2+o0.1i/0.1i+w2c -O -K >> $ps
gmt pstext -R0/10/0/15 -JX10/10 -X0.5c -Y12.5c -N -O \
    -F+f12p,0,black+jLB >> $ps << EOF
1.0 3.0 EGM96: 15 arc minute resolution grid based on the gravitational force of the Earth
EOF
gmt psconvert Geoid_Ker.ps -A1.6c -E720 -Tj -Z
\end{lstlisting}

%%%%%%%%%%%%%%%%%%%%%%%%%%%%%%%%%%%%%%%%%%
\begin{adjustwidth}{-\extralength}{0cm}
%\printendnotes[custom] % Un-comment to print a list of endnotes

\reftitle{References}
%\nocite{*}

%=====================================
% References, variant A: external bibliography
%=====================================
\bibliography{Kerguelen.bib}

%%%%%%%%%%%%%%%%%%%%%%%%%%%%%%%%%%%%%%%%%%
\end{adjustwidth}
\end{document}
